\chapter{Bronchiectasis}

\section*{Definition and Overview}
Bronchiectasis is a long-term lung condition in which the airways (bronchi) become permanently widened, damaged, and scarred. Because the damaged airways can no longer clear mucus effectively, secretions accumulate within them, leading to a chronic productive cough, recurrent chest infections, and progressive damage. The structural changes are irreversible, but modern management can control symptoms, reduce the frequency and severity of flare-ups, and slow further deterioration. Bronchiectasis is distinct from asthma and chronic obstructive pulmonary disease (COPD), although the three conditions frequently coexist. It is a chronic illness that is best managed jointly by a respiratory specialist and a general practitioner, with physiotherapy as a central pillar of care.

\section*{Epidemiology}
Bronchiectasis is less common than asthma or COPD but is increasingly recognised, partly because high-resolution CT scanning detects the condition more readily than in the past. It can affect any age group but is most often diagnosed in middle-aged and older adults, with a slight female predominance in many populations. Higher-risk groups include people who suffered severe childhood infections, those with primary antibody deficiencies or autoimmune disease, and people with cystic fibrosis. In Australia, bronchiectasis is significantly more prevalent in some Aboriginal and Torres Strait Islander communities, particularly among children. Smoking, chronic respiratory infection, and recurrent aspiration also contribute to disease development.

\section*{Aetiology and Pathophysiology}
\begin{itemize}
    \item \textbf{Severe or repeated infections:} whooping cough, measles, pneumonia, tuberculosis, and other severe respiratory infections, especially in childhood, can permanently damage the airways
    \item \textbf{Cystic fibrosis:} a genetic disorder producing thick, tenacious mucus that injures the airways from early life
    \item \textbf{Immune dysfunction:} antibody deficiencies and autoimmune conditions such as rheumatoid arthritis predispose to airway damage
    \item \textbf{Aspiration:} repeated inhalation of food, fluid, or stomach acid into the lungs
    \item \textbf{Mechanical obstruction:} a foreign body, tumour, or enlarged lymph node blocking an airway and trapping secretions
    \item \textbf{Mechanism:} chronic inflammation destroys the airway wall, including the cilia and mucus-secreting glands that normally clear secretions; the airway widens and becomes floppy, mucus pools within it, bacteria flourish, and repeated infection causes further injury in a self-perpetuating cycle
\end{itemize}

\section*{Clinical Features}
\begin{itemize}
    \item A chronic productive cough, often present for months or years
    \item Coughing up large volumes of phlegm that may be thick, discoloured, or foul-smelling
    \item Recurrent chest infections and exacerbations that temporarily worsen symptoms
    \item Shortness of breath, worse with exertion
    \item Wheeze or a coarse, rattling chest sound
    \item Fatigue and a general feeling of being unwell
    \item Finger clubbing (widening and rounding of the fingertips) in some people
    \item Coughing up blood (haemoptysis), usually small amounts but occasionally large
    \item Poor weight gain or weight loss in children and in severe adult disease
\end{itemize}

\section*{Differential Diagnosis}
\begin{itemize}
    \item \textbf{COPD:} chronic cough and sputum production in a smoker, with airflow obstruction on testing
    \item \textbf{Asthma:} wheeze and cough, but typically without the large sputum volumes and permanent airway damage of bronchiectasis
    \item \textbf{Chronic bronchitis:} a persistent productive cough, usually smoking-related
    \item \textbf{Tuberculosis:} chronic infection with cough, fever, night sweats, and weight loss
    \item \textbf{Lung cancer:} a new or changing cough, haemoptysis, or unexplained weight loss
    \item \textbf{Heart failure:} pulmonary oedema causing breathlessness and cough
    \item \textbf{Allergic bronchopulmonary aspergillosis:} an allergic response to the fungus Aspergillus that produces similar airway damage
\end{itemize}

\section*{When to Seek Medical Care}
See a doctor if a cough lasts more than three weeks, if phlegm is produced on most days, if chest infections occur repeatedly, or if any amount of blood is coughed up. People with known bronchiectasis should seek early review for changes in sputum, fever, or increased breathlessness, since prompt antibiotic treatment limits flare-ups.

\textbf{Call 000 (or local emergency services) for:}
\begin{itemize}
    \item Coughing up a large amount of blood, or bleeding that does not stop
    \item Sudden or severe difficulty breathing
    \item Chest pain that is severe or lasts longer than a few minutes
    \item A very high fever or a feeling of being extremely unwell during a flare-up
    \item Blue lips, confusion, or extreme drowsiness
\end{itemize}

\section*{Diagnosis and Investigations}
\begin{itemize}
    \item \textbf{History:} the pattern and duration of cough and sputum, childhood infections, smoking, family history, and underlying conditions such as rheumatoid arthritis or immunodeficiency
    \item \textbf{Examination:} chest auscultation, oxygen saturation measurement, and inspection for finger clubbing
    \item \textbf{Chest imaging:} high-resolution CT of the chest is the diagnostic test of choice and demonstrates the characteristic dilated airways; a chest X-ray may be normal
    \item \textbf{Sputum microbiology:} culture of phlegm to identify the organisms colonising the airway and their antibiotic sensitivities
    \item \textbf{Lung function tests:} spirometry and gas transfer measurements to quantify the degree of airflow obstruction and functional impairment
    \item \textbf{Blood tests:} full blood count, inflammatory markers, immunoglobulin levels, and autoantibody screening where indicated
    \item \textbf{Specialised tests:} sweat chloride and genetic testing for cystic fibrosis, allergy testing for aspergillosis, and bronchoscopy in selected cases
\end{itemize}

\section*{Management}
\begin{itemize}
    \item \textbf{Airway clearance physiotherapy:} daily positioning, breathing techniques, and devices such as positive expiratory pressure (PEP) valves to loosen and drain mucus; this is the cornerstone of treatment
    \item \textbf{Bronchodilators and inhaled steroids:} inhalers to open the airways, especially where asthma or COPD coexist
    \item \textbf{Antibiotics:} targeted courses for exacerbations based on sputum cultures; some people receive regular low-dose or inhaled antibiotics to reduce infections
    \item \textbf{Mucus-thinning therapy:} nebulised hypertonic saline or mucolytic agents to make sputum easier to clear
    \item \textbf{Vaccination:} annual influenza and pneumococcal vaccination to reduce infection risk
    \item \textbf{Treatment of the underlying cause:} immunoglobulin replacement for antibody deficiency, disease-modifying therapy for rheumatoid arthritis, or measures to reduce aspiration
    \item \textbf{Management of haemoptysis:} small bleeds usually settle with rest; major bleeding may require bronchial artery embolisation
    \item \textbf{Surgery:} rarely, resection of a severely damaged lung segment in selected patients
    \item \textbf{Pulmonary rehabilitation:} supervised exercise and education to improve fitness, symptoms, and quality of life
\end{itemize}

\section*{Complications}
\begin{itemize}
    \item Frequent chest infections requiring repeated or prolonged antibiotic courses
    \item Progressive decline in lung function with worsening breathlessness over the years
    \item Haemoptysis, which can occasionally be life-threatening
    \item Cor pulmonale, a form of right-heart strain caused by advanced lung disease
    \item Respiratory failure, requiring long-term oxygen therapy in advanced disease
    \item Antibiotic resistance and side effects from long-term treatment
    \item Chronic fatigue, weight loss, and impaired quality of life
\end{itemize}

\section*{Prognosis and Outcomes}
Bronchiectasis is a chronic, incurable condition, and the structural damage to the airways is permanent. With consistent treatment, however, many people remain stable and active for many years, and the goals of management are symptom control, fewer exacerbations, and slower progression. The outlook depends on the underlying cause, the extent of lung involvement, and adherence to treatment. People with severe or widespread disease, frequent infections, or an underlying condition such as cystic fibrosis tend to experience a greater decline in lung function. Early diagnosis and rigorous daily airway clearance give the best chance of preserving lung health.

\section*{Prevention and Self-Care}
\begin{itemize}
    \item Do not smoke, and avoid second-hand smoke and other lung irritants
    \item Perform prescribed airway clearance exercises every day, even when feeling well
    \item Keep influenza and pneumococcal vaccinations up to date
    \item Complete the full course of any prescribed antibiotics
    \item Drink plenty of fluids to keep mucus thinner and easier to clear
    \item Stay physically active within your limits, and consider pulmonary rehabilitation
    \item Practise good hand hygiene and avoid close contact with people with respiratory infections
    \item Attend regular reviews with your respiratory specialist and general practitioner
\end{itemize}

\section*{Sources and Further Reading}
The following organisations provide reliable information on bronchiectasis for patients, students, and clinicians.

\begin{itemize}
    \item NHS. \textit{Bronchiectasis: symptoms, diagnosis, and treatment.} https://www.nhs.uk/conditions/bronchiectasis/
    \item Mayo Clinic. \textit{Bronchiectasis: overview, causes, and management.} https://www.mayoclinic.org/
    \item MedlinePlus. \textit{Bronchiectasis health information.} https://medlineplus.gov/bronchiectasis.html
    \item MSD Manual. \textit{Bronchiectasis: professional overview and management.} https://www.msdmanuals.com/
    \item National Institute for Health and Care Excellence. \textit{Bronchiectasis (non-cystic fibrosis) guidance.} https://www.nice.org.uk/
    \item Thoracic Society of Australia and New Zealand. \textit{Adult and paediatric bronchiectasis guidelines and resources.} https://www.thoracic.org.au/
\end{itemize}
